% Generated draft: requires mathematical and source review.
\documentclass[11pt]{article}
\usepackage[margin=25mm]{geometry}
\usepackage{fontspec}
\setmainfont{Latin Modern Roman}
\usepackage{amsmath,amssymb,mathtools,unicode-math}
\setmathfont{Latin Modern Math}
\usepackage{enumitem,needspace,xurl,hyperref}
\setlength{\parindent}{0pt}
\setlength{\parskip}{5pt}
\setlength{\emergencystretch}{3em}
\newcommand{\N}{\mathbb{N}}
\newcommand{\Z}{\mathbb{Z}}
\newcommand{\Q}{\mathbb{Q}}
\newcommand{\R}{\mathbb{R}}
\newcommand{\C}{\mathbb{C}}
\newcommand{\F}{\mathbb{F}}
\newcommand{\A}{\mathbb{A}}
\newcommand{\PP}{\mathbb{P}}
\newcommand{\E}{\mathbb{E}}
\newcommand{\eps}{\varepsilon}
\newcommand{\dd}{\,\mathrm{d}}
\begin{document}
\section*{UnsolvedMath ID: 30004024}
\textbf{Source problem code:} catalogue entry
\subsection{Definitions and mathematical statement}
For an integer \(n\geq 1\), a convex body in \(\R^n\) means a compact convex set with nonempty interior. For any such body \(L\), let \(|L|\) denote its \(n\)-dimensional Lebesgue measure, and define its barycenter and covariance matrix with respect to normalized Lebesgue measure by
\[
b_L=\frac{1}{|L|}\int_L x\,\dd x,
\qquad
\operatorname{Cov}(L)=\frac{1}{|L|}\int_L (x-b_L)(x-b_L)^{\mathsf T}\,\dd x.
\]
If \(b_K=0\), then the origin lies in the interior of \(K\), and its polar convex body is
\[
K^\circ=\{y\in\R^n:\langle x,y\rangle\leq 1\text{ for every }x\in K\}.
\]
The covariance of \(K^\circ\) is centered at its own barycenter \(b_{K^\circ}\); no assumption that this barycenter vanishes is imposed.

The supplied determinant question is: does there exist a constant \(C>0\), independent of both \(n\) and \(K\), such that for every integer \(n\geq 1\) and every convex body \(K\subseteq\R^n\) with \(b_K=0\),
\[
\det\!\left(\operatorname{Cov}(K)\operatorname{Cov}(K^\circ)\right)
\leq \left(\frac{C}{n}\right)^{2n}?
\]
The supplied dated curation describes this determinant relation as established while labeling the overall record partially solved and referring to broader, only partially resolved trace-based duality relations. No precise trace-based question is supplied, so none is added here. These status descriptions are supplied assessments, not independently verified mathematical conclusions.
\subsection{Short English statement}
Is there an absolute constant \(C>0\) such that, in every dimension \(n\geq 1\), the determinant of the product of the covariance matrices of a convex body with barycenter zero and its polar is at most \((C/n)^{2n}\)? Each covariance uses the uniform probability measure and is centered at that body's own barycenter. The supplied curation calls this relation established, but that assessment requires independent verification.
\subsection{Sources}
\begin{itemize}
\item Supplied local UnsolvedMath registry record, permanent ID \(30004024\), source code \texttt{catalogue entry}, title ``Covariance Determinant Bounds for Polar Convex Bodies''; project registry \texttt{data/unsolvedmath/problems.json}, with \texttt{catalog/problems.json} designated as the supplementary registry. No individual proposer is identified in the supplied record.
\item Supplied report citation: \emph{Convex Geometry and its Applications} (2019), Oberwolfach Reports, pp.~3217--3281, DOI \texttt{10.4171/owr/2018/54}. \url{https://doi.org/10.4171/owr/2018/54}. This is also the supplied literature reference labeled ``Oberwolfach report 54.'' The record assigns proposed year 2018; the relationship between these year labels has not been independently checked.
\item Supplied literature reference: ``Klartag, isotropic constants and Mahler volumes.'' \url{https://arxiv.org/abs/1710.08084}.
\item Supplied literature reference: ``Klartag--Lehec.'' \url{https://doi.org/10.1007/s00039-025-00718-w}.
\item Literature assessment dated exactly \texttt{2026-08-21}: supplied dated curation, not independently verified. It labels the record ``partially solved'' and states that the main relation is established, linking isotropic constants and polarity through the covariance determinant and the slicing inequality, described there as affirmed by Klartag--Lehec. It also states that broader trace-based duality relations remain only partial. The supplied background expressly requests independent verification of time-sensitive status and cited claims.
\end{itemize}
\end{document}
